\pdfbookmark{Общая характеристика работы}{characteristic}             % Закладка pdf
\section*{Общая характеристика работы}

\newcommand{\actuality}{\pdfbookmark[1]{Актуальность}{actuality}\underline{\textbf{\actualityTXT}}}
\newcommand{\progress}{\pdfbookmark[1]{Разработанность темы}{progress}\underline{\textbf{\progressTXT}}}
\newcommand{\aim}{\pdfbookmark[1]{Цели}{aim}\underline{{\textbf\aimTXT}}}
\newcommand{\tasks}{\pdfbookmark[1]{Задачи}{tasks}\underline{\textbf{\tasksTXT}}}
\newcommand{\aimtasks}{\pdfbookmark[1]{Цели и задачи}{aimtasks}\aimtasksTXT}
\newcommand{\object}{\pdfbookmark[1]{Объект исследования}{object}\underline{\textbf{\objectTXT}}}
\newcommand{\subject}{\pdfbookmark[1]{Предмет исследования}{subject}\underline{\textbf{\subjectTXT}}}
\newcommand{\novelty}{\pdfbookmark[1]{Научная новизна}{novelty}\underline{\textbf{\noveltyTXT}}}
\newcommand{\tinfluence}{\pdfbookmark[1]{Теоретическая значимость}{tinfluence}\underline{\textbf{\tinfluenceTXT}}}
\newcommand{\pinfluence}{\pdfbookmark[1]{Практическая значимость}{pinfluence}\underline{\textbf{\pinfluenceTXT}}}
\newcommand{\methods}{\pdfbookmark[1]{Методология и методы исследования}{methods}\underline{\textbf{\methodsTXT}}}
\newcommand{\defpositions}{\pdfbookmark[1]{Положения, выносимые на защиту}{defpositions}\underline{\textbf{\defpositionsTXT}}}
\newcommand{\passport}{\pdfbookmark[1]{Соответствие паспорту специальности}{passport}\underline{\textbf{\passportTXT}}}
\newcommand{\reliability}{\pdfbookmark[1]{Достоверность}{reliability}\underline{\textbf{\reliabilityTXT}}}
\newcommand{\probation}{\pdfbookmark[1]{Апробация}{probation}\underline{\textbf{\probationTXT}}}
\newcommand{\contribution}{\pdfbookmark[1]{Личный вклад}{contribution}\underline{\textbf{\contributionTXT}}}
\newcommand{\publications}{\pdfbookmark[1]{Публикации}{publications}\underline{\textbf{\publicationsTXT}}}

\input{common/characteristic} % Характеристика работы по структуре во введении и в автореферате не отличается (ГОСТ Р 7.0.11, пункты 5.3.1 и 9.2.1), потому её загружаем из одного и того же внешнего файла, предварительно задав форму выделения некоторым параметрам

%Диссертационная работа была выполнена при поддержке грантов \dots

%\underline{\textbf{Объем и структура работы.}} Диссертация состоит из~введения,
%четырех глав, заключения и~приложения. Полный объем диссертации
%\textbf{ХХХ}~страниц текста с~\textbf{ХХ}~рисунками и~5~таблицами. Список
%литературы содержит \textbf{ХХX}~наименование.

\pdfbookmark{Содержание работы}{description}                          % Закладка pdf
\section*{Содержание работы}
Во \underline{\textbf{введении}} обосновывается актуальность
исследований, проводимых в~рамках данной диссертационной работы,
приводится обзор научной литературы по~изучаемой проблеме,
формулируется цель, ставятся задачи работы, излагается научная новизна
и практическая значимость представляемой работы.

\underline{\textbf{Первая глава}} носит постановочный характер и посвящена анализу проблемы определения достаточного размера обучающей выборки для параметрических моделей машинного обучения. Объектом называется пара $(\mathbf{x}, y)$, где $\mathbf{x} \in \mathbb{X} \subseteq \mathbb{R}^n$ "--- вектор признакового описания, $y \in \mathbb{Y}$ "--- целевая переменная. В задаче регрессии $\mathbb{Y} = \mathbb{R}$, в задаче классификации $\mathbb{Y} = \{1, \ldots, K\}$. Выборка размера~$m$ задается множеством $\mathfrak{D}_m = \{(\mathbf{x}_i, y_i)\}_{i=1}^m$ из $m$ независимых реализаций пар~$(\mathbf{x}, y)$, распределенных согласно неизвестному распределению $p(\mathbf{x}, y)$.

Параметрической моделью называется семейство функций $f_{\mathbf{w}}: \mathbb{X} \to \mathbb{Y}$, задаваемое вектором параметров $\mathbf{w} \in \mathbb{W} \subset \mathbb{R}^N$. Выбор модели из этого семейства проводится с помощью эмпирической функции потерь, вычисляемой на заданной выборке:
\begin{equation}
\label{eq:synopsis-empirical-loss}
\mathcal{L}_m(\mathbf{w}) = \frac{1}{m} \sum_{i=1}^m \ell(f_{\mathbf{w}}(\mathbf{x}_i), y_i),
\end{equation}
где $\ell: \mathbb{Y} \times \mathbb{Y} \to \mathbb{R}^+$ "--- функция потерь на одном объекте (например, квадратичная или кросс-энтропийная). Обучение модели таким образом есть решение оптимизационной задачи
\begin{equation}
    \label{eq:synopsis-min-problem}
    \mathbf{w}_m^* \in \argmin_{\mathbf{w} \in \mathbb{W}} \mathcal{L}_m(\mathbf{w}).
\end{equation}

Далее рассмотрены основные существующие подходы к оценке достаточного размера выборки "--- статистические, байесовские, информационно-теоретические и основанные на оценках сложности модели "--- и показаны их ограничения применительно к современным многопараметрическим нейросетевым архитектурам. На этой основе обоснована необходимость построения новых критериев, опирающихся не на внешние характеристики качества или асимптотические статистические гипотезы, а на локальные свойства самой оптимизационной задачи.

В качестве ключевого объекта исследования предложено использовать оптимизационную поверхность функции потерь и ее локальную геометрию в окрестности точки минимума. Сформулирована основная идея диссертации: рассматривать стабилизацию ландшафта функции потерь при увеличении объема обучающей выборки как математическую основу для построения критериев достаточного размера выборки и интерпретации законов масштабирования параметрических моделей.

В заключение главы сформулированы задачи диссертационного исследования и описана логика последующих глав: анализ спектральных свойств матрицы Гессе, построение критериев стабилизации ландшафта, вывод критериев достаточного размера выборки и исследование вероятностных линейных моделей как частного случая параметрических моделей.

\underline{\textbf{Во второй главе}} исследуются спектральные свойства матрицы Гессе функции потерь нейронной сети и выводятся архитектурно-зависимые оценки локальной кривизны оптимизационной поверхности.

В качестве основного объекта рассматривается матрица Гессе эмпирической функции потерь
\begin{equation}
\mathbf{H}^{(m)}(\mathbf{w})=\nabla_{\mathbf w}^2 \mathcal L_m(\mathbf w).
\end{equation}
Для анализа спектрально значимой части второй производной используется разложение Гаусса--Ньютона
\begin{equation}
\mathbf H(\mathbf w)=\mathbf G(\mathbf w)+\mathbf R(\mathbf w),\qquad
\mathbf G(\mathbf w)=\mathbf J(\mathbf w)^\top \mathbf A(\mathbf w)\mathbf J(\mathbf w),
\end{equation}
где $\mathbf J(\mathbf w)=\partial \mathbf z/\partial \mathbf w$ "--- якобиан выхода модели по параметрам, а $\mathbf A(\mathbf w)=\nabla_{\mathbf z}^2\ell(\mathbf z,\mathbf y)$ "--- матрица кривизны функции потерь по выходу модели.

Для широкого класса матрично-представимых нейронных сетей в работе получена блочная факторизация гаусс--ньютоновской компоненты. Если $\mathbf w=\mathrm{col}(\mathbf w^{(1)},\dots,\mathbf w^{(L)})$ "--- разбиение параметров по слоям, а
\[
\mathbf S^{(p)}(\mathbf w)=\frac{\partial \mathbf z}{\partial \mathbf z^{(p)}},\qquad
\mathbf P^{(p)}(\mathbf w)=\frac{\partial \mathbf z^{(p)}}{\partial \mathbf w^{(p)}},
\]
то для $(p,q)$-го блока матрицы $\mathbf G(\mathbf w)$ доказано представление
\begin{equation}
\mathbf G^{(p,q)}(\mathbf w)=
\mathbf P^{(p)}(\mathbf w)^\top
\mathbf S^{(p)}(\mathbf w)^\top
\mathbf A(\mathbf w)
\mathbf S^{(q)}(\mathbf w)
\mathbf P^{(q)}(\mathbf w).
\end{equation}
Полученная факторизация, проиллюстрированная на рисунке~\ref{fig:synopsis-gn-factorization}, позволяет разделить вклад архитектуры сети, функции потерь и параметризации слоев в локальную кривизну.

\begin{figure}[ht]
    \centering
    \resizebox{\linewidth}{!}{\begin{tikzpicture}[
    >=Latex,
    font=\small,
    every node/.style={align=center},
    blockmat/.style={draw, line width=0.5pt},
    flow/.style={-Latex, thick},
    factor/.style={
        draw,
        rounded corners=2pt,
        minimum height=8mm,
        minimum width=14mm,
        inner sep=2pt
    }
]

% -------------------------------------------------
% Left: big block matrix G
% -------------------------------------------------
\def\W{4.8}
\def\H{4.8}

\coordinate (Gsw) at (0,0);
\coordinate (Gne) at (\W,\H);

\draw[blockmat] (Gsw) rectangle (Gne);

% grid 4x4
\foreach \x in {1.2,2.4,3.6}
    \draw[blockmat] (\x,0) -- (\x,\H);
\foreach \y in {1.2,2.4,3.6}
    \draw[blockmat] (0,\y) -- (\W,\y);

% labels inside matrix
\node at (2.4,5.25) {$\mathbf{G}(\mathbf{w})=\big[\mathbf{G}^{(p, q)}(\mathbf{w})\big]_{p, q=1}^{L}$};

\node at (0.6,4.2) {$\mathbf{G}^{(1, 1)}$};
\node at (1.8,4.2) {$\mathbf{G}^{(1, 2)}$};
\node at (3.0,4.2) {$\cdots$};
\node at (4.2,4.2) {$\mathbf{G}^{(1, L)}$};

\node at (0.6,3.0) {$\mathbf{G}^{(2, 1)}$};
\node at (1.8,3.0) {$\mathbf{G}^{(2, 2)}$};
\node at (3.0,3.0) {$\cdots$};
\node at (4.2,3.0) {$\mathbf{G}^{(2, L)}$};

\node at (0.6,1.8) {$\vdots$};
\node at (1.8,1.8) {$\vdots$};
\node at (3.0,1.8) {$\ddots$};
\node at (4.2,1.8) {$\vdots$};

\node at (0.6,0.6) {$\mathbf{G}^{(L, 1)}$};
\node at (1.8,0.6) {$\mathbf{G}^{(L, 2)}$};
\node at (3.0,0.6) {$\cdots$};
\node at (4.2,0.6) {$\mathbf{G}^{(L, L)}$};

% row/col annotations
\node[left=4mm] at (0,2.4) {\rotatebox{90}{\small параметры $p$-го слоя}};
\node[below=4mm] at (2.4,0) {\small параметры $q$-го слоя};

% single arrow from selected block
\coordinate (selected) at (3.0,1.8);

% -------------------------------------------------
% Right: factorization of one block
% -------------------------------------------------
\node[factor] at (8.0,1.8) (Pt) {$\mathbf{P}^{(p)\top}$};
\node[factor, right=0.18cm of Pt] (St) {$\mathbf{S}^{(p)\top}$};
\node[factor, right=0.18cm of St] (A) {$\mathbf{A}(\mathbf{w})$};
\node[factor, right=0.18cm of A] (S) {$\mathbf{S}^{(q)}$};
\node[factor, right=0.18cm of S] (P) {$\mathbf{P}^{(q)}$};

\draw[flow] (selected) -- (Pt.west);

\node[above=6mm of A] {\small факторизация отдельного блока};
\node[below=6mm of A] {$\displaystyle
\mathbf{G}^{(p,q)}(\mathbf{w})
=
\mathbf{P}^{(p)\top}
\mathbf{S}^{(p)\top}
\mathbf{A}(\mathbf{w})
\mathbf{S}^{(q)}
\mathbf{P}^{(q)}
$};

\end{tikzpicture}}
    \caption{Блочная факторизация гаусс--ньютоновской компоненты матрицы Гессе для матрично-представимой нейронной сети}
    \label{fig:synopsis-gn-factorization}
\end{figure}

На основе этого результата для $L$-слойной полносвязной сети получена верхняя оценка спектральной нормы гаусс--ньютоновской компоненты
\begin{equation}
\|\mathbf G(\mathbf w)\|_2
\le
L M_{\mathbf x}^2 M_{\mathbf A}
\left(M_{\mathbf W}M_{\sigma}\right)^{2L-2},
\end{equation}
где $M_{\mathbf x}$ ограничивает норму входа, $M_{\mathbf A}$ "--- спектральную норму матрицы $\mathbf A(\mathbf w)$, $M_{\mathbf W}$ "--- нормы весовых матриц, а $M_{\sigma}$ "--- нормы матриц производных активаций.

Для сверточной нейронной сети доказана аналогичная оценка
\begin{equation}
\|\mathbf G(\mathbf w)\|_2
\le
L M_{\mathbf x}^2 M_{\mathbf A} K_{\max}
\left(M_{\mathcal W}M_{\sigma}\right)^{2L-2},
\end{equation}
где $M_{\mathcal W}$ ограничивает $\ell_1$-нормы сверточных ядер, а множитель $K_{\max}$ определяется максимальным размером рецептивного поля слоя. Тем самым показано, что для сверточных сетей верхняя оценка локальной кривизны зависит от локальности связей и разделения весов.

Кроме того, для полносвязной сети с активацией ReLU и кросс-энтропийной функцией потерь получено частное следствие
\begin{equation}
\|\mathbf G(\mathbf w)\|_2
\le
\sqrt{2}\,L M_{\mathbf x}^2 M_{\mathbf W}^{2L-2}.
\end{equation}

Проведенные вычислительные эксперименты для полносвязных и сверточных архитектур на задаче классификации показывают, что на поздних этапах обучения выполняется соотношение
\[
\|\mathbf G(\mathbf w)\|_2 \approx \|\mathbf H(\mathbf w)\|_2,
\]
что подтверждает возможность использования гаусс--ньютоновской компоненты как практически значимой характеристики локальной кривизны оптимизационной поверхности.

\underline{\textbf{В третьей главе}} исследуется стабилизация ландшафта функции потерь нейронной сети при увеличении объема обучающей выборки. В отличие от классического анализа, основанного на сходимости эмпирического риска как скалярной величины, в главе рассматривается изменение локальной геометрии оптимизационной поверхности в окрестности точки минимума.

В качестве исходного объекта изучается разность $\mathcal L_{k+1}(\mathbf w)-\mathcal L_k(\mathbf w)$ эмпирических функций потерь на выборках размеров $k$ и $k+1$. В окрестности точки локального минимума $\mathbf w_k^*$ функции $\mathcal L_k(\mathbf w)$ эта разность аппроксимируется формулой второго порядка
\begin{multline}
\mathcal L_{k+1}(\mathbf w)-\mathcal L_k(\mathbf w)
=
\mathcal L_{k+1}(\mathbf w_k^*)-\mathcal L_k(\mathbf w_k^*)
+
\mathbf g^{(k+1)}(\mathbf w_k^*)^\top(\mathbf w-\mathbf w_k^*)
+
\\
+
\frac12
(\mathbf w-\mathbf w_k^*)^\top
\left(
\mathbf H^{(k+1)}(\mathbf w_k^*)-\mathbf H^{(k)}(\mathbf w_k^*)
\right)
(\mathbf w-\mathbf w_k^*),
\end{multline}
где $\mathbf g^{(k+1)}(\mathbf w_k^*)$ "--- градиент функции потерь на выборке из $k+1$ объектов, а $\mathbf H^{(k)}(\mathbf w_k^*)$ и $\mathbf H^{(k+1)}(\mathbf w_k^*)$ "--- соответствующие матрицы Гессе. Тем самым изменение ландшафта описывается тремя составляющими: сдвигом значения функции потерь в точке минимума, градиентным рассогласованием (см. иллюстрацию на рисунке~\ref{fig:synopsis-losses-assumption}) и изменением локальной кривизны.

\begin{figure}[ht]
    \centering
    \includegraphics[width=0.62\linewidth]{losses_assumption.pdf}
    \caption{Иллюстрация условия слабого градиентного рассогласования в окрестности локального минимума}
    \label{fig:synopsis-losses-assumption}
\end{figure}

На этой основе вводится общий критерий $p$-сходимости ландшафта функции потерь:
\begin{equation}
\Delta_p(k+1)
=
\int
\left|
\mathcal L_{k+1}(\mathbf w)-\mathcal L_k(\mathbf w)
\right|^p
q(\mathbf w)\,d\mathbf w,
\end{equation}
где функция $q(\mathbf w)$ задает предпочтение по выбору точек в пространстве параметров. Такой критерий позволяет единообразно рассматривать как изменение функции потерь в одной фиксированной точке, так и усредненные характеристики стабилизации поверхности в полном пространстве параметров или в выделенном подпространстве, как это проиллюстрировано на рисунке~\ref{fig:synopsis-landscape-criteria}.

\begin{figure}[ht]
    \centering
    \begin{minipage}{0.32\linewidth}
        \centering
        \includegraphics[width=\linewidth]{criterion_abs.pdf}
        \small а)
    \end{minipage}\hfill
    \begin{minipage}{0.32\linewidth}
        \centering
        \includegraphics[width=\linewidth]{criterion_squared.pdf}
        \small б)
    \end{minipage}\hfill
    \begin{minipage}{0.32\linewidth}
        \centering
        \includegraphics[width=\linewidth]{criterion_squared_subspace.pdf}
        \small в)
    \end{minipage}
    \caption{Иллюстрация введенных критериев стабилизации ландшафта функции потерь: а) абсолютный одноточечный критерий; б) среднеквадратичный критерий в полном пространстве параметров; в) среднеквадратичный критерий в подпространстве наибольшей кривизны}
    \label{fig:synopsis-landscape-criteria}
\end{figure}

Для абсолютного одноточечного критерия
\begin{equation}
\Delta_1(k+1)
=
\left|
\mathcal L_{k+1}(\mathbf w)-\mathcal L_k(\mathbf w)
\right|,
\qquad
\mathbf w\in\mathcal U_R(\mathbf w_k^*),
\end{equation}
получена оценка
\begin{equation}
\Delta_1(k+1)
\le
\frac{2M_\ell+M_{\mathbf g}R+M_{\mathbf H}R^2}{k+1}
=
\mathcal O(k^{-1}),
\end{equation}
где $M_\ell$, $M_{\mathbf g}$ и $M_{\mathbf H}$ ограничивают соответственно значения функции потерь на одном объекте, градиента и матрицы Гессе в окрестности минимума.

Для среднеквадратичного критерия
\begin{equation}
\Delta_2(k+1)
=
\int
\big(
\mathcal L_{k+1}(\mathbf w)-\mathcal L_k(\mathbf w)
\big)^2
q(\mathbf w)\,d\mathbf w
\end{equation}
при выборе гауссовского распределения
\begin{equation}
q(\mathbf w)=\mathcal N(\mathbf w_k^*,\sigma^2\mathbf I_N)
\end{equation}
доказана оценка
\begin{equation}
\Delta_2(k+1)
\le
\frac{
12M_\ell^2
+
3\sigma^2M_{\mathbf g}^2
+
3\sigma^4(N^2+2N)M_{\mathbf H}^2
}{(k+1)^2}
=
\mathcal O(k^{-2}).
\end{equation}
Таким образом, после усреднения квадрата изменения ландшафта по распределению точек получается более быстрая скорость убывания по сравнению с одноточечным критерием.

Далее вводится подпространство наибольшей кривизны
\begin{equation}
\mathcal S_D=\operatorname{Im}(\mathbf U_D),
\qquad
\mathbf U_D=[\mathbf u_1,\dots,\mathbf u_D],
\end{equation}
натянутое на собственные векторы $\mathbf u_1,\dots,\mathbf u_D$, соответствующие $D$ наибольшим собственным значениям матрицы Гессе $\mathbf H^{(k)}(\mathbf w_k^*)$. В этом подпространстве рассматривается критерий
\begin{equation}
\Delta_2^{(D)}(k+1)
=
\int
\big(
\mathcal L_{k+1}(\mathbf w)-\mathcal L_k(\mathbf w)
\big)^2
q(\mathbf w)\,d\mathbf w,
\qquad
\mathbf w\in \mathbf w_k^*+\mathcal S_D.
\end{equation}
При гауссовском распределении в координатах подпространства доказано, что
\begin{equation}
\Delta_2^{(D)}(k+1)
\le
\frac{
12M_\ell^2
+
3\sigma^2M_{\mathbf g}^2
+
3\sigma^4(D^2+2D)M_{\mathbf H}^2
}{(k+1)^2}
=
\mathcal O(k^{-2}).
\end{equation}
Тем самым показано, что ограничение анализа подпространством главных собственных направлений матрицы Гессе сохраняет тот же порядок сходимости, что и в полном пространстве параметров, но позволяет заменить размерность $N$ на существенно меньшую размерность $D$.

В частном случае отсутствия нулевого и линейного вкладов критерий в подпространстве наибольшей кривизны выражается непосредственно через приращения ведущих собственных значений матрицы Гессе:
\begin{equation}
\Delta_{2}^{(D)}(k+1)
=
\frac{\sigma^4}{4}
\left(
2\sum_{i=1}^{D}
\left(
\lambda_i^{(k+1)}-\lambda_i^{(k)}
\right)^2
+
\left(
\sum_{i=1}^{D}
\left(
\lambda_i^{(k+1)}-\lambda_i^{(k)}
\right)
\right)^2
\right).
\end{equation}
Это дает спектральную интерпретацию стабилизации ландшафта функции потерь.

Проведенные вычислительные эксперименты для полносвязной нейронной сети на задаче классификации MNIST подтверждают полученные теоретические результаты. Среднеквадратичный критерий в полном пространстве и подпространственные критерии демонстрируют убывание, близкое к закону $\mathcal O(k^{-2})$, тогда как абсолютный одноточечный критерий убывает медленнее и обладает большей вариативностью. Дополнительные визуализации спектра матрицы Гессе и двумерных сечений поверхности функции потерь показывают, что основные изменения ландшафта действительно концентрируются в направлениях наибольшей локальной кривизны.

\underline{\textbf{В четвертой главе}} результаты, полученные ранее для спектральных свойств матрицы Гессе и для критериев стабилизации ландшафта функции потерь, объединяются в единую схему, позволяющую построить архитектурно-зависимые критерии достаточного размера выборки для параметрических моделей машинного обучения.

Исходным пунктом является разложение пообъектной матрицы Гессе в форме Гаусса--Ньютона
\begin{equation}
\mathbf H_i(\mathbf w)=\mathbf G_i(\mathbf w)+\mathbf R_i(\mathbf w),
\end{equation}
где $\mathbf G_i(\mathbf w)$ есть гаусс--ньютоновская компонента, а $\mathbf R_i(\mathbf w)$ "--- остаточный член. Для фиксированной архитектуры $\mathcal A$ вводится локальная архитектурно-зависимая константа кривизны
\begin{equation}
M_{\mathbf G}(\mathcal A)
=
\sup_{\mathbf w\in\mathcal U_R(\mathbf w^*)}
\sup_i
\|\mathbf G_i(\mathbf w)\|_2,
\end{equation}
а также величина
\begin{equation}
\delta_R
=
\sup_{\mathbf w\in\mathcal U_R(\mathbf w^*)}
\sup_i
\|\mathbf R_i(\mathbf w)\|_2,
\end{equation}
характеризующая вклад остаточного члена. Тогда для любой точки локальной окрестности минимума справедливо неравенство
\begin{equation}
\|\mathbf H_i(\mathbf w)\|_2
\le
M_{\mathbf G}(\mathcal A)+\delta_R.
\end{equation}
Это позволяет во всех критериях стабилизации ландшафта заменить общую константу $M_{\mathbf H}$ на величину $M_{\mathbf G}(\mathcal A)+\delta_R$, зависящую от архитектуры модели.

На этой основе достаточный размер выборки определяется через выбранный критерий стабилизации ландшафта функции потерь:
\begin{equation}
m^*
=
\inf\{m\in\mathbb N:\Delta(m)\le\varepsilon\},
\end{equation}
где $\varepsilon>0$ "--- заранее заданный уровень точности. Тем самым выборка считается достаточной, если при дальнейшем увеличении числа обучающих объектов изменение локальной геометрии функции потерь становится несущественным по выбранному критерию.

Для абсолютного одноточечного критерия получена оценка достаточного размера выборки вида
\begin{equation}
m_1^*(\varepsilon,\mathcal A)
=
\frac{
2M_{\ell}
+
M_{\mathbf g}R
+
\left(M_{\mathbf G}(\mathcal A)+\delta_R\right)R^2
}{\varepsilon},
\end{equation}
где $M_{\ell}$ ограничивает значения функции потерь на отдельных объектах, $M_{\mathbf g}$ "--- величину градиентного рассогласования, а $R$ "--- радиус локальной окрестности минимума.

Для среднеквадратичного критерия в полном пространстве параметров при гауссовской функции предпочтения
\begin{equation}
q(\mathbf w)=\mathcal N(\mathbf w^*,\sigma^2\mathbf I_N)
\end{equation}
получена оценка
\begin{equation}
m_2^*(\varepsilon,\mathcal A)
=
\left[
\frac{
12M_{\ell}^2
+
3\sigma^2M_{\mathbf g}^2
+
3\sigma^4(N^2+2N)\left(M_{\mathbf G}(\mathcal A)+\delta_R\right)^2
}{\varepsilon}
\right]^{1/2}.
\end{equation}
Аналогично, для среднеквадратичного критерия в подпространстве наибольшей кривизны размерности $D$ доказана оценка
\begin{equation}
m_3^*(\varepsilon,\mathcal A,D)
=
\left[
\frac{
12M_{\ell}^2
+
3\sigma^2M_{\mathbf g}^2
+
3\sigma^4(D^2+2D)\left(M_{\mathbf G}(\mathcal A)+\delta_R\right)^2
}{\varepsilon}
\right]^{1/2}.
\end{equation}
По сравнению с полным пространством параметров здесь размерностный множитель зависит не от полной размерности модели $N$, а только от размерности выбранного информативного подпространства $D$.

Подстановка полученных во второй главе архитектурно-зависимых оценок для $M_{\mathbf G}(\mathcal A)$ дает явные критерии достаточного размера выборки для полносвязных и сверточных нейронных сетей. Для $L$-слойной полносвязной сети без смещений локальная константа кривизны имеет вид
\begin{equation}
M_{\mathbf G}^{\mathrm{FC}}
=
L M_{\mathbf x}^2 M_{\mathbf A}(M_{\mathbf W}M_{\sigma})^{2L-2},
\end{equation}
а для $L$-слойной сверточной сети без смещений
\begin{equation}
M_{\mathbf G}^{\mathrm{CNN}}
=
L M_{\mathbf x}^2 M_{\mathbf A}K_{\max}(M_{\mathcal W}M_{\sigma})^{2L-2}.
\end{equation}
Тем самым достаточный размер выборки выражается через глубину архитектуры, нормы весовых операторов, свойства функций активации и, в сверточном случае, размер рецептивного поля.

Полученные формулы показывают, что при прочих равных увеличение глубины приводит к росту необходимого объема данных для стабилизации ландшафта функции потерь. В полносвязных сетях этот рост связан с накоплением кривизны в длинных произведениях плотных операторов, тогда как в сверточных сетях существенную роль играет локальность связей, отражаемая множителем $K_{\max}$. В результате архитектура модели интерпретируется не только как класс реализуемых функций, но и как фактор, определяющий скорость стабилизации локальной геометрии задачи оптимизации.

Помимо строгих критериев достаточного размера выборки, в главе предложена более общая интерпретация полученных результатов в терминах законов масштабирования параметрических моделей. Для этого рассматривается декомпозиция значения функции потерь в обученной точке $\mathbf w_m^*$:
\begin{equation}
\mathcal L_m(\mathbf w_m^*)
=
\underbrace{\mathcal L(\mathbf w_{\infty}^*)}_{E}
+
\underbrace{\bigl(\mathcal L(\mathbf w_m^*)-\mathcal L(\mathbf w_{\infty}^*)\bigr)}_{\text{ошибка модели}}
+
\underbrace{\bigl(\mathcal L_m(\mathbf w_m^*)-\mathcal L(\mathbf w_m^*)\bigr)}_{\mathcal E(m)}.
\end{equation}
Здесь $E$ есть неустранимая ошибка задачи, среднее слагаемое отражает ограниченность класса моделей, а величина $\mathcal E(m)$ описывает вклад конечности обучающей выборки.

Для ошибки выразимости вводится функция
\begin{equation}
\Phi(N)
=
\inf_{\mathbf w\in\mathcal W_N}\mathcal L(\mathbf w)-\mathcal L(\mathbf w_{\infty}^*),
\end{equation}
характеризующая ограниченность класса моделей с числом параметров $N$. Далее формулируется структурное предположение о том, что в локальном статистическом режиме вклад конечности выборки наследует геометрический масштаб стабилизации ландшафта:
\begin{equation}
|\mathcal E(m)|
\le
\frac{C\left(M_{\mathbf G}(\mathcal A)+\delta_R\right)}{m^{\rho}},
\qquad
\rho\in(0,1].
\end{equation}
При этом получается схема локального закона масштабирования
\begin{equation}
\mathcal L_m(\mathbf w_m^*)
\le
E
+
\Phi(N)
+
\frac{C\left(M_{\mathbf G}(\mathcal A)+\delta_R\right)}{m^{\rho}}.
\end{equation}
В режиме, когда остаточный член мал по сравнению с гаусс--ньютоновской компонентой, эта формула принимает упрощенный вид
\begin{equation}
\mathcal L_m(\mathbf w_m^*)
\lesssim
E
+
\Phi(N)
+
\frac{CM_{\mathbf G}(\mathcal A)}{m^{\rho}}.
\end{equation}

При дополнительном предположении о степенном убывании ошибки выразимости
\begin{equation}
\Phi(N)\sim \frac{A}{N^{\alpha}}
\end{equation}
получается выражение
\begin{equation}
\mathcal L_m(\mathbf w_m^*)
\lesssim
E
+
\frac{A}{N^{\alpha}}
+
\frac{CM_{\mathbf G}(\mathcal A)}{m^{\rho}},
\end{equation}
в котором одно слагаемое отвечает за ограниченность класса моделей, а второе "--- за конечность обучающей выборки и зависит от локальной геометрии оптимизационной поверхности.

Далее рассматривается вычислительно-ограниченный режим обучения, в котором общий бюджет пропорционален произведению числа параметров модели и объема обучающих данных:
\begin{equation}
\mathcal C\propto Nm.
\end{equation}
Если в фиксированном режиме масштабирования архитектуры выполнено
\begin{equation}
M_{\mathbf G}(N)\sim BN^{\gamma},
\end{equation}
то оптимальное соотношение между размером модели и размером выборки удовлетворяет асимптотикам
\begin{equation}
N^*
\propto
\mathcal C^{\frac{\rho}{\alpha+\gamma+\rho}},
\qquad
m^*
\propto
\mathcal C^{\frac{\alpha+\gamma}{\alpha+\gamma+\rho}}.
\end{equation}
Тем самым оптимальный баланс между моделью и данными определяется не только скоростью убывания ошибки выразимости, но и тем, насколько быстро растет локальная кривизна архитектуры в выбранном режиме масштабирования.

Проведенные вычислительные эксперименты на полносвязных и сверточных архитектурах для задачи классификации MNIST подтверждают основные качественные выводы главы. Во-первых, критерий стабилизации ландшафта $\Delta_2(m)$ убывает с ростом объема выборки по закону, согласующемуся с теоретической оценкой порядка $\mathcal O(m^{-2})$. Во-вторых, внутри каждого семейства архитектур достаточный размер выборки возрастает с увеличением глубины. В-третьих, структурная аппроксимация обобщающего разрыва степенным законом по $m$ показывает, что вклад конечности обучающей выборки действительно допускает представление в форме data term с архитектурно-зависимыми параметрами.

\underline{\textbf{В пятой главе}} рассматриваются критерии достаточного размера выборки для вероятностных линейных моделей. В отличие от предыдущих глав, где достаточность выборки исследовалась через стабилизацию ландшафта функции потерь и его локальной геометрии, здесь изучается вероятностная постановка, в которой основным объектом анализа становится стабилизация оценок параметров и соответствующих апостериорных распределений.

Рассматривается вероятностная модель линейной регрессии
\begin{equation}
p(y | \mathbf x,\mathbf w)=\mathcal N(y | \mathbf w^\top \mathbf x,\sigma^2),
\end{equation}
где $\mathbf x\in\mathbb R^n$ "--- вектор признаков, $y\in\mathbb R$ "--- отклик, $\mathbf w\in\mathbb R^n$ "--- вектор параметров, а $\sigma^2>0$ "--- дисперсия шума. Для выборки $\mathfrak D_k=(\mathbf X_k,\mathbf y_k)$ оценка максимума правдоподобия определяется как решение задачи наименьших квадратов
\begin{equation}
\hat{\mathbf w}_k
=
\argmin_{\mathbf w\in\mathbb R^n}
\|\mathbf y_k-\mathbf X_k\mathbf w\|_2^2.
\end{equation}
При гауссовском априорном распределении параметров апостериорное распределение имеет сопряженный вид
\begin{equation}
p(\mathbf w | \mathfrak D_k)=\mathcal N(\mathbf w | \mathbf m_k,\boldsymbol\Sigma_k),
\end{equation}
где параметры $\mathbf m_k$ и $\boldsymbol\Sigma_k$ явно выражаются через $\mathbf X_k$ и $\mathbf y_k$.

На этой основе в главе вводятся четыре критерия достаточного размера выборки, основанные на различных характеристиках стабилизации решения. Первые два критерия, проиллюстрированные на рисунке~\ref{fig:likelihood}, используют поведение функции правдоподобия в точке оценки максимума правдоподобия. D-критерий основан на уменьшении дисперсии правдоподобия при варьировании подвыборок, а M-критерий "--- на уменьшении изменения его математического ожидания при переходе от выборки размера $k$ к выборке размера $k+1$.

\begin{figure}[ht]
    \centering
    \includegraphics[width=0.5\linewidth]{likelihood.pdf}
    \caption{Визуализация критериев D-достаточности и M-достаточности размера выборки в вероятностной модели линейной регрессии}
    \label{fig:likelihood}
\end{figure}

Оставшиеся два критерия, проиллюстрированные на рисунке~\ref{fig:posterior}, используют близость апостериорных распределений параметров на схожих подвыборках. KL-критерий определяется через дивергенцию Кульбака--Лейблера
\begin{equation}
D_{\mathrm{KL}}(p_k\|p_{k+1})
=
\int p_k(\mathbf w)\log\frac{p_k(\mathbf w)}{p_{k+1}(\mathbf w)}\,d\mathbf w,
\end{equation}
а S-критерий "--- через функцию близости s-score, нормированную так, что для совпадающих распределений она равна единице. Тем самым достаточный размер выборки определяется через стабилизацию либо правдоподобия, либо полного апостериорного распределения параметров.

\begin{figure}[ht]
    \centering
    \includegraphics[width=0.88\linewidth]{posterior_ru.pdf}
    \caption{Визуализация критериев KL-достаточности и S-достаточности размера выборки в вероятностной модели линейной регрессии}
    \label{fig:posterior}
\end{figure}

В главе формализуется понятие корректности критерия достаточности: критерий называется корректным, если соответствующая функция сходится к своему предельному значению при увеличении объема выборки. Для M-, KL- и S-критериев доказываются теоремы корректности, связывающие ее со сходимостью моментов апостериорного распределения. В частности, показывается, что если
\begin{equation}
\|\mathbf m_{k+1}-\mathbf m_k\|_2\to 0,
\qquad
\|\boldsymbol\Sigma_{k+1}-\boldsymbol\Sigma_k\|_F\to 0
\quad\text{при}\quad k\to\infty,
\end{equation}
то M- и KL-критерии являются корректными, а для S-критерия достаточно сходимости апостериорных средних.

Далее для вероятностной линейной регрессии с гауссовским априором доказывается теорема, дающая достаточные условия для выполнения этих сходимостей. Предполагается ограниченность признаков и отклика, а также условие роста минимального собственного значения матрицы Грама:
\begin{equation}
\lambda_{\min}(\mathbf X_k^\top \mathbf X_k)=\omega(\sqrt{k}).
\end{equation}
При этих предположениях устанавливается, что
\begin{equation}
\|\mathbf m_{k+1}-\mathbf m_k\|_2\to 0,
\qquad
\|\boldsymbol\Sigma_{k+1}-\boldsymbol\Sigma_k\|_F\to 0
\quad\text{при}\quad k\to\infty.
\end{equation}
Тем самым получаются проверяемые условия, обеспечивающие корректность введенных критериев достаточности в модели линейной регрессии.

Помимо стабилизации самих распределений параметров, в главе исследуется стабилизация полного байесовского прогноза на отложенной тестовой выборке. Пусть
\begin{equation}
p(\mathbf y_{\mathrm{test}} | \mathbf X_{\mathrm{test}},\mathbf X_k,\mathbf y_k)
\end{equation}
обозначает предиктивное распределение, полученное по обучающей подвыборке размера $k$. Доказывается, что при тех же условиях на данные и матрицу признаков
\begin{equation}
D_{\mathrm{KL}}
\left(
p(\mathbf y_{\mathrm{test}} | \mathbf X_{\mathrm{test}},\mathbf X_k,\mathbf y_k)
\big\|
p(\mathbf y_{\mathrm{test}} | \mathbf X_{\mathrm{test}},\mathbf X_{k+1},\mathbf y_{k+1})
\right)
\to 0
\quad\text{при}\quad k\to\infty.
\end{equation}
Это означает, что увеличение объема обучающей выборки приводит не только к стабилизации апостериорного распределения параметров, но и к стабилизации байесовского прогноза на новых данных.

Проведенные вычислительные эксперименты включают как синтетические данные, сгенерированные из модели линейной регрессии, так и реальные регрессионные выборки из открытых источников. В экспериментах строятся последовательности апостериорных распределений на схожих подвыборках, исследуются функции $D(k)$, $M(k)$, $KL(k)$ и $S(k)$, а также анализируется зависимость достаточного размера выборки от порогового параметра $\varepsilon$. Результаты показывают, что при увеличении размера подвыборки значения функций $D(k)$, $M(k)$ и $KL(k)$ стремятся к нулю, а значение функции $S(k)$ стремится к единице. Это согласуется с доказанными в главе теоретическими утверждениями.

Таким образом, в содержательной части автореферата последовательно раскрывается общая логика диссертационного исследования. Сначала формулируется задача определения достаточного размера обучающей выборки для параметрических моделей, затем исследуются архитектурно-зависимые свойства локальной кривизны оптимизационной поверхности, после чего строятся критерии стабилизации ландшафта функции потерь и выводятся критерии достаточного размера выборки. Отдельно рассматривается вероятностная линейная регрессия как частный случай параметрических моделей, для которого достаточность выборки описывается через стабилизацию распределения оптимальных параметров и байесовского прогноза.

\FloatBarrier
\pdfbookmark{Заключение}{conclusion}                                  % Закладка pdf
В \underline{\textbf{заключении}} приведены основные результаты работы, которые заключаются в следующем:
%% Согласно ГОСТ Р 7.0.11-2011:
%% 5.3.3 В заключении диссертации излагают итоги выполненного исследования, рекомендации, перспективы дальнейшей разработки темы.
%% 9.2.3 В заключении автореферата диссертации излагают итоги данного исследования, рекомендации и перспективы дальнейшей разработки темы.
\begin{enumerate}
    \item Получены верхние оценки спектральной нормы гаусс"--~ньютоновской компоненты матрицы Гессе для полносвязных и сверточных нейронных сетей.
    \item Предложены критерии стабилизации ландшафта функции потерь при увеличении объема обучающей выборки.
    \item Для введенных критериев получены оценки скорости сходимости в окрестности точки минимума.
    \item Построены критерии достаточного размера выборки для параметрических моделей на основе стабилизации оптимизационной поверхности.
    \item Предложена теоретическая схема интерпретации связи между локальной геометрией функции потерь, архитектурой модели и законами масштабирования.
    \item Для вероятностных линейных моделей построены критерии достаточного размера выборки, основанные на стабилизации апостериорного распределения параметров.
    \item Вычислительные эксперименты качественно подтвердили согласованность полученных теоретических результатов с наблюдаемым поведением моделей.
\end{enumerate}


\pdfbookmark{Литература}{bibliography}                                % Закладка pdf

\ifdefmacro{\microtypesetup}{\microtypesetup{protrusion=false}}{} % не рекомендуется применять пакет микротипографики к автоматически генерируемому списку литературы
\urlstyle{rm}                               % ссылки URL обычным шрифтом
\ifnumequal{\value{bibliosel}}{0}{% Встроенная реализация с загрузкой файла через движок bibtex8
    \renewcommand{\bibname}{\large \bibtitleauthor}
    \nocite{*}
    \insertbiblioauthor           % Подключаем Bib-базы
    %\insertbiblioexternal   % !!! bibtex не умеет работать с несколькими библиографиями !!!
}{% Реализация пакетом biblatex через движок biber
    % Цитирования.
    %  * Порядок перечисления определяет порядок в библиографии (только внутри подраздела, если `\insertbiblioauthorgrouped`).
    %  * Если не соблюдать порядок "как для \printbibliography", нумерация в `\insertbiblioauthor` будет кривой.
    %  * Если цитировать каждый источник отдельной командой "--- найти некоторые ошибки будет проще.
    %
    % \nocite{}
    \nocite{kiselev2025ssdlb}
    \nocite{kiselev2025ssdpdp}
    \nocite{kiselev2024unraveling}
    \nocite{kiselev2024conf66mipt}
    \nocite{meshkov2024convnets}
    \nocite{kiselev2025conf67mipt}
    \nocite{meshkov2025conf67mipt}
    \nocite{kiselev2025mmro}

    \ifnumgreater{\value{usefootcite}}{0}{
        \begin{refcontext}[labelprefix={}]
            \ifnum \value{bibgrouped}>0
                \insertbiblioauthorgrouped    % Вывод всех работ автора, сгруппированных по источникам
            \else
                \insertbiblioauthor      % Вывод всех работ автора
            \fi
        \end{refcontext}
    }{
        \ifnum \totvalue{citeexternal}>0
            \begin{refcontext}[labelprefix=A]
                \ifnum \value{bibgrouped}>0
                    \insertbiblioauthorgrouped    % Вывод всех работ автора, сгруппированных по источникам
                \else
                    \insertbiblioauthor      % Вывод всех работ автора
                \fi
            \end{refcontext}
        \else
            \ifnum \value{bibgrouped}>0
                \insertbiblioauthorgrouped    % Вывод всех работ автора, сгруппированных по источникам
            \else
                \insertbiblioauthor      % Вывод всех работ автора
            \fi
        \fi
         % \insertbiblioauthorimportant  % Вывод наиболее значимых работ автора (определяется в файле characteristic во второй section)
        \begin{refcontext}[labelprefix={}]
            \insertbiblioexternal            % Вывод списка литературы, на которую ссылались в тексте автореферата
        \end{refcontext}
        % Невидимый библиографический список для подсчёта количества внешних публикаций
        % Используется, чтобы убрать приставку "А" у работ автора, если в автореферате нет
        % цитирований внешних источников.
        \printbibliography[heading=nobibheading, section=0, env=countexternal, keyword=biblioexternal, resetnumbers=true]%
    }
}
\ifdefmacro{\microtypesetup}{\microtypesetup{protrusion=true}}{}
\urlstyle{tt}                               % возвращаем установки шрифта ссылок URL
